% Options for packages loaded elsewhere
\PassOptionsToPackage{unicode}{hyperref}
\PassOptionsToPackage{hyphens}{url}
\PassOptionsToPackage{dvipsnames,svgnames,x11names}{xcolor}
%
\documentclass[
  letterpaper,
  DIV=11,
  numbers=noendperiod]{scrartcl}

\usepackage{amsmath,amssymb}
\usepackage{iftex}
\ifPDFTeX
  \usepackage[T1]{fontenc}
  \usepackage[utf8]{inputenc}
  \usepackage{textcomp} % provide euro and other symbols
\else % if luatex or xetex
  \usepackage{unicode-math}
  \defaultfontfeatures{Scale=MatchLowercase}
  \defaultfontfeatures[\rmfamily]{Ligatures=TeX,Scale=1}
\fi
\usepackage{lmodern}
\ifPDFTeX\else  
    % xetex/luatex font selection
\fi
% Use upquote if available, for straight quotes in verbatim environments
\IfFileExists{upquote.sty}{\usepackage{upquote}}{}
\IfFileExists{microtype.sty}{% use microtype if available
  \usepackage[]{microtype}
  \UseMicrotypeSet[protrusion]{basicmath} % disable protrusion for tt fonts
}{}
\makeatletter
\@ifundefined{KOMAClassName}{% if non-KOMA class
  \IfFileExists{parskip.sty}{%
    \usepackage{parskip}
  }{% else
    \setlength{\parindent}{0pt}
    \setlength{\parskip}{6pt plus 2pt minus 1pt}}
}{% if KOMA class
  \KOMAoptions{parskip=half}}
\makeatother
\usepackage{xcolor}
\setlength{\emergencystretch}{3em} % prevent overfull lines
\setcounter{secnumdepth}{-\maxdimen} % remove section numbering
% Make \paragraph and \subparagraph free-standing
\makeatletter
\ifx\paragraph\undefined\else
  \let\oldparagraph\paragraph
  \renewcommand{\paragraph}{
    \@ifstar
      \xxxParagraphStar
      \xxxParagraphNoStar
  }
  \newcommand{\xxxParagraphStar}[1]{\oldparagraph*{#1}\mbox{}}
  \newcommand{\xxxParagraphNoStar}[1]{\oldparagraph{#1}\mbox{}}
\fi
\ifx\subparagraph\undefined\else
  \let\oldsubparagraph\subparagraph
  \renewcommand{\subparagraph}{
    \@ifstar
      \xxxSubParagraphStar
      \xxxSubParagraphNoStar
  }
  \newcommand{\xxxSubParagraphStar}[1]{\oldsubparagraph*{#1}\mbox{}}
  \newcommand{\xxxSubParagraphNoStar}[1]{\oldsubparagraph{#1}\mbox{}}
\fi
\makeatother


\providecommand{\tightlist}{%
  \setlength{\itemsep}{0pt}\setlength{\parskip}{0pt}}\usepackage{longtable,booktabs,array}
\usepackage{calc} % for calculating minipage widths
% Correct order of tables after \paragraph or \subparagraph
\usepackage{etoolbox}
\makeatletter
\patchcmd\longtable{\par}{\if@noskipsec\mbox{}\fi\par}{}{}
\makeatother
% Allow footnotes in longtable head/foot
\IfFileExists{footnotehyper.sty}{\usepackage{footnotehyper}}{\usepackage{footnote}}
\makesavenoteenv{longtable}
\usepackage{graphicx}
\makeatletter
\newsavebox\pandoc@box
\newcommand*\pandocbounded[1]{% scales image to fit in text height/width
  \sbox\pandoc@box{#1}%
  \Gscale@div\@tempa{\textheight}{\dimexpr\ht\pandoc@box+\dp\pandoc@box\relax}%
  \Gscale@div\@tempb{\linewidth}{\wd\pandoc@box}%
  \ifdim\@tempb\p@<\@tempa\p@\let\@tempa\@tempb\fi% select the smaller of both
  \ifdim\@tempa\p@<\p@\scalebox{\@tempa}{\usebox\pandoc@box}%
  \else\usebox{\pandoc@box}%
  \fi%
}
% Set default figure placement to htbp
\def\fps@figure{htbp}
\makeatother

\usepackage{booktabs}
\usepackage{longtable}
\usepackage{array}
\usepackage{multirow}
\usepackage{wrapfig}
\usepackage{float}
\usepackage{colortbl}
\usepackage{pdflscape}
\usepackage{tabu}
\usepackage{threeparttable}
\usepackage{threeparttablex}
\usepackage[normalem]{ulem}
\usepackage{makecell}
\usepackage{xcolor}
\KOMAoption{captions}{tableheading}
\makeatletter
\@ifpackageloaded{caption}{}{\usepackage{caption}}
\AtBeginDocument{%
\ifdefined\contentsname
  \renewcommand*\contentsname{Tabla de contenidos}
\else
  \newcommand\contentsname{Tabla de contenidos}
\fi
\ifdefined\listfigurename
  \renewcommand*\listfigurename{Listado de Figuras}
\else
  \newcommand\listfigurename{Listado de Figuras}
\fi
\ifdefined\listtablename
  \renewcommand*\listtablename{Listado de Tablas}
\else
  \newcommand\listtablename{Listado de Tablas}
\fi
\ifdefined\figurename
  \renewcommand*\figurename{Figura}
\else
  \newcommand\figurename{Figura}
\fi
\ifdefined\tablename
  \renewcommand*\tablename{Tabla}
\else
  \newcommand\tablename{Tabla}
\fi
}
\@ifpackageloaded{float}{}{\usepackage{float}}
\floatstyle{ruled}
\@ifundefined{c@chapter}{\newfloat{codelisting}{h}{lop}}{\newfloat{codelisting}{h}{lop}[chapter]}
\floatname{codelisting}{Listado}
\newcommand*\listoflistings{\listof{codelisting}{Listado de Listados}}
\makeatother
\makeatletter
\makeatother
\makeatletter
\@ifpackageloaded{caption}{}{\usepackage{caption}}
\@ifpackageloaded{subcaption}{}{\usepackage{subcaption}}
\makeatother

\ifLuaTeX
\usepackage[bidi=basic]{babel}
\else
\usepackage[bidi=default]{babel}
\fi
\babelprovide[main,import]{spanish}
% get rid of language-specific shorthands (see #6817):
\let\LanguageShortHands\languageshorthands
\def\languageshorthands#1{}
\usepackage{bookmark}

\IfFileExists{xurl.sty}{\usepackage{xurl}}{} % add URL line breaks if available
\urlstyle{same} % disable monospaced font for URLs
\hypersetup{
  pdftitle={Bitácora 1},
  pdflang={es},
  colorlinks=true,
  linkcolor={blue},
  filecolor={Maroon},
  citecolor={Blue},
  urlcolor={Blue},
  pdfcreator={LaTeX via pandoc}}


\title{Bitácora 1}
\author{}
\date{}

\begin{document}
\maketitle


\begin{center}

\textbf{Universidad de Costa Rica}


\vspace{0.5cm}

\textbf{Escuela de Matemática}

\textbf{Departamento de Matemática y Ciencias Actuariales}

\vspace{0.5cm}

\textbf{Proyecto de Investigación}

\textbf{Análisis del cambio de la probabilidad de decremento ${}_{n}q_{x}^{(c)}$  para grupos etarios por  causas de muerte de acuerdo a la clasificación ICD-10 en países de Centroamérica [2015, 2018].}

\vspace{0.5cm}

\textbf{Integrantes}

\textbf{Sebastián Miranda Ramírez  C4H274 }

\textbf{Benjamín Gutiérrez Padua  C4F813}

\textbf{Kevin David Calderón Martínez  C4D511}

\textbf{Gabriel de Jesús Chaves Esquivel  C4E273}

\vspace{0.5cm}

\textbf{Profesor}

\textbf{Maikol Solís Chacón} 

\vspace{0.5cm}

\textbf{2026}


\end{center}

\newpage

\section{1. Planificación}\label{planificaciuxf3n}

\subsection{1.1 Pregunta de
investigación}\label{pregunta-de-investigaciuxf3n}

¿Qué patrones de variabilidad se pueden estimar en la probabilidad de
decremento \({}_{n}q_{x}^{(c)}\) por categoría de causas de la ICD-10 en
Países Centroamericanos en los años 2015-2018?

\subsubsection{1.1.1 Definición de la
idea}\label{definiciuxf3n-de-la-idea}

La Organización Mundial de la Salud (WHO, por sus siglas en inglés) ha
desempeñado históricamente un papel central en la estandarización
internacional de las estadísticas de mortalidad y de las causas de
muerte. En este marco, los países remiten datos provenientes de sus
sistemas de registro civil y la OMS los compila y armoniza conforme a
estándares internacionales, en particular los vinculados a la
Clasificación Internacional de Enfermedades (ICD, por sus siglas en
inglés) y a las Regulaciones de Nomenclatura de 1967. (World Health
Organization {[}WHO{]}, 2026).

De esta forma, la idea inicial del proyecto consiste en \textbf{analizar
cómo cambia la probabilidad de decremento por categorías de causas
seleccionadas según grupos de edad en países de Centroamérica entre
2015-2018, utilizando la WHO Mortality Database y trabajando con las
categorías de la ICD-10.} En particular, el interés se centra en
estudiar el objeto actuarial \({}_{n}q_{x}^{(c)}\), entendido como la
probabilidad de que una persona de edad \(x\) muera en el intervalo
\((x, x+n]\) debido a la causa \(c\), en presencia de otras causas
competidoras. Para este fin, se utilizarán tanto el enfoque clásico como
el enfoque bayesiano, con el propósito de analizar si ambos producen
estimaciones semejantes o si comienzan a diferir.

La elección de este enfoque también responde a antecedentes empíricos.
En América Latina, la literatura reciente, en particular Calazans y
Queiroz (2020), sugiere que, aunque existen patrones comunes en la
mortalidad adulta por causa de muerte, sus cambios a lo largo del tiempo
no son idénticos en todos los países. No obstante, estos estudios no
analizan la evolución de una edad a otra, sino la mortalidad adulta en
un rango agregado de edades. Esto vuelve razonable seguir estudiando el
fenómeno con mayor detalle y desde un punto de vista de evolución de una
edad a otra, lo que vuelve aún más adecuado el objeto
\({}_{n}q_{x}^{(c)}\).

\subsubsection{1.1.2 Conceptualización de la
idea}\label{conceptualizaciuxf3n-de-la-idea}

La idea a conceptualizar es: analizar cómo cambia la probabilidad de
decremento por categorías de causas seleccionadas de la ICD-10 según
grupos etarios en países de Centroamérica, utilizando la WHO Mortality
Database y trabajando con las categorías de la ICD-10.

Según el Diccionario de la lengua española, analizar significa someter
algo a un análisis, y análisis se entiende como un estudio detallado de
algo. Segundo, probabilidad se define como cálculo o determinación
cuantitativa de la posibilidad de que se verifique un suceso. Decremento
significa, en su uso general, disminución; causa es aquello que se
considera fundamento u origen de algo; edad es el tiempo que ha vivido
una persona y también cada uno de los períodos en que se divide la vida;
grupo es una pluralidad de seres o cosas que forman un conjunto;
categoría es cada una de las clases establecidas dentro de una
actividad; y determinado alude a algo particular o específico.

Asimismo, categorías de causas seleccionadas de la ICD-10 significa, en
términos prácticos, usar las categorías amplias definidas por la WHO,
como enfermedades infecciosas y parasitarias \textbf{(I)} o tumores
\textbf{(II)} que son bastante más manejables que la lista detallada de
tres o cuatro caracteres. Por otro lado, según grupos de edad es una
frase importante, porque la base de la OMS no entrega una cohorte
individual seguida en el tiempo, sino conteos agregados por edades o
grupos de edad, con formatos etarios que pueden variar entre países y
años. Finalmente, utilizar la WHO Mortality Database se refiere a
trabajar con los datos agregados por país proporcionados de manera
estandarizada por la WHO.

Con algunos conceptos llevados al contexto del proyecto, es importante
diferenciar lo siguiente: que probabilidad se entenderá como una
magnitud entre 0 y 1 que expresa la proporción esperada de veces que
ocurrirá un evento. Decremento no se usará en el sentido coloquial de
simple disminución, sino en su sentido actuarial; la salida de una vida
del estado inicial por una causa específica, basado en la terminología
de la Society of Actuaries, la cual define tablas de múltiples
decrementos como aquellas que registran vidas, número de decrementos y
probabilidades de decremento y supervivencia para más de una causa. En
este caso, los decrementos de interés serán las muertes por causas
mutuamente excluyentes.

\subsubsection{1.1.3 Identificación de
tensiones}\label{identificaciuxf3n-de-tensiones}

\subsubsection{1.1.3 Identificación}\label{identificaciuxf3n}

Una de las primeras tensiones que podría tener este tema es el uso de la
base de datos. Al ser una base tan amplia y con gran cantidad de datos
el manejo de la misma puede de salirse fácilmente de control si no se
consideran los cuidados necesarios, es por esto que es vital el reducir
los datos solamente a los fundamentalmente requeridos y formular el
problema sobre un subconjunto de la base.

La segunda tensión es actuarial. La cantidad \({}_{n}q_{x}^{(c)}\) no
viene observada directamente en la base de la OMS. Lo que la base ofrece
son conteos de muertes por causa y edad, y poblaciones de referencia por
edad. Por tanto, el proyecto tendrá que decidir cómo pasar de esos
conteos a una estimación de decremento por causas.

La tercera tensión viene directamente del uso y comparación del enfoque
estadístico clásico y bayesiano, puesto que hay que asegurarse que ambos
se apliquen al mismo parámetro y conjunto de datos. Además, al momento
de la comparación surgen problemas como el calcular la función de error
de cada uno de los dos enfoques con efectividad y sus intervalos de
confianza, ya que, de lo contrario podría surgir una comparación poco
precisa dando resultados erróneos.

\subsubsection{1.1.4 Reformulación de la idea en modo de
pregunta}\label{reformulaciuxf3n-de-la-idea-en-modo-de-pregunta}

1 ¿Cómo \textbf{estimar} los elementos de \textbf{patología de base} que
inciden en la \textbf{intensidad} de la \textbf{mortalidad} en la
\textbf{población} de ciertos \textbf{Países Centroamericanos} tomando
de referencia los datos de 2015-2018 de la OMS?~

2 ¿Qué \textbf{patrones de variabilidad} se pueden \textbf{estimar} en
la \textbf{probabilidad de decremento} nqx(c) por c\textbf{ategoría de
causas} de la ICD-10 en \textbf{Países Centroamericanos} en los años
2015-2018?

3 ¿Cómo se \textbf{distribuyen} las \textbf{muertes} por
\textbf{categorías de causas} de la ICD-10 según \textbf{grupos de
etario}s en \textbf{Países de Centroamérica} entre 2015 y 2018?

4 ¿Cómo puede \textbf{estimarse} desde un \textbf{enfoque frecuentista y
uno bayesiano}, una tasa de \textbf{mortalidad} por causas seleccionadas
en \textbf{distintos estratos de la base}?

Se determinó que la pregunta 2 encapsula de modo más integral los
objetivos que se desean alcanzar en la consumación del proyecto,
mientras que las otras tres preguntas se descartan como eje principal
por razones de alcance y pertinencia. La primera es útil para la parte
conceptual del agregado de categorías de causas de muerte, pero por sí
sola no alcanza a estudiar el cambio de la probabilidad de decremento.
La tercera cuenta con un enfoque más descriptivo y se centra en el uso
de herramientas probabilísticas que suponen más un medio que un fin con
los objetivos de la investigación.La cuarta, por su parte, sí aborda de
manera directa la comparación entre el enfoque clásico y el enfoque
bayesiano, pero convierte lo convierte en un tema principal que, dentro
del proyecto, funciona más bien como apoyo metodológico que como
objetivo sustantivo central. Por ello, la segunda pregunta resulta la
más viable y coherente con los datos disponibles, con el objeto de
estudio escogido y con el interés de investigación ya planteado.

\subsubsection{1.1.5 Argumentación de la
pregunta}\label{argumentaciuxf3n-de-la-pregunta}

1 ¿Cómo \textbf{estimar} los elementos de \textbf{patología de base} que
inciden en la \textbf{intensidad} de la \textbf{mortalidad} en la
\textbf{población} de ciertos \textbf{Países Centroamericanos} tomando
de referencia los datos de 2015-2018 de la OMS?~

\textbf{Contraargumentos}

\begin{itemize}
\item
  \textbf{Lógica:} Se podría cuestionar que el delimitar a cierto
  intervalo temporal inciden negativamente en la calidad de la
  generalización sobre las observaciones realizadas.
\item
  \textbf{Ética:} Partir de factores patológicos dados no canaliza la
  espontaneidad de eventos que repercuten en la variabilidad de causas
  de muerte posibles.
\item
  \textbf{Emocional:} Los factores patológicos más significativos de los
  que se ha demostrado su incidencia en la mortalidad para personas
  supone un componente previamente investigado, por lo que el enfoque no
  aporta al público una perspectiva de estudio novedosa.
\end{itemize}

\textbf{Argumentos}

\begin{itemize}
\item
  \textbf{Lógica:} El modelar elementos que inciden en la intensidad de
  la mortalidad para ciertas clasificaciones de causas de perecimiento
  permite reconocer de qué manera se ve intensificada la mortalidad por
  características del grupo delimitado.
\item
  \textbf{Ética:} Se puede establecer un marco de referencia para el
  análisis para la evolución de la incidencia de ciertas enfermedades en
  diferentes edades de las personas.
\item
  \textbf{Emocional:} Se genera material para el análisis para el
  estudio y toma de decisiones por autoridades competentes para la
  prevención de determinados factores que suscitan el perecimiento de
  personas.
\end{itemize}

2 ¿Qué \textbf{patrones de variabilidad} se pueden \textbf{estimar} en
la \textbf{probabilidad de decremento} nqx(c) por c\textbf{ategoría de
causas} de la ICD-10 en \textbf{Países Centroamericanos} en los años
2015-2018?

\textbf{Contraargumentos}

\begin{itemize}
\item
  \textbf{Lógica:} La probabilidad de decrementos puede ser un objeto
  matemáticamente difícil de construir si no se modela adecuadamente la
  cadena de Markov para diferentes estados.
\item
  \textbf{Ética:} El identificar patrones de cambio entre intervalos de
  edad puede involucrar la interpretación subjetiva del investigador.
\item
  \textbf{Emocional:} La facultad de apertura e interpretabilidad pueden
  verse afectadas para personas no especializadas en los temas
  abordados.~
\end{itemize}

\textbf{Argumentos}

\begin{itemize}
\item
  \textbf{Lógica:} Se estima con un enfoque sólido teórico el
  comportamiento de la evolución de la mortalidad para diferentes
  intervalos de edad.
\item
  \textbf{Ética: ~F}aculta la descripción de la causa de mortalidad
  entre grupos de etarios.
\item
  \textbf{Emocional:} El enfoque le proporciona a personas
  especializadas una gama de datos, en específico las probabilidades de
  decremento, para entender cómo evolucionan las características para
  diferentes grupos etarios de determinadas causas de mortalidad.
\end{itemize}

3 ¿Cómo se \textbf{distribuyen} las \textbf{muertes} por
\textbf{categorías de causas} de la ICD-10 según \textbf{grupos de
etario}s en \textbf{Países de Centroamérica} entre 2015 y 2018?

\textbf{Contraargumentos}

\begin{itemize}
\item
  \textbf{Lógica:} Al indexar una diversidad de factores de patología de
  base dentro de un solo mancomunado conjunto elimina la heterogeneidad
  de los datos y puede~ asignar proporciones de muertes no realistas
  para ciertas patologías estudiadas por individual.
\item
  \textbf{Ética: S}e pierde la distribución marginal característica para
  cada respectiva causa de muerte.
\item
  \textbf{Emocional:} La construcción del agregado de factores de muerte
  y su justificación puede estimular una opacidad al momento de
  interpretar los resultados para personas no especializadas.~
\end{itemize}

\textbf{Argumentos}

\begin{itemize}
\item
  \textbf{Lógica:} Faculta la descripción intuitiva de qué elementos
  inciden mayoritariamente en mortalidad para diferentes estratos de
  edad.
\item
  \textbf{Ética:} Promueve el análisis de determinadas categorías de
  causa de muerte para población en diferentes grupos de etarios.~
\item
  \textbf{Emocional:} Proporciona una herramienta intuitiva para que
  personas no necesariamente especializadas o capacitadas en menesteres
  técnicos comprendan factores que inciden en la mortalidad en distintos
  intervalos de edad.
\end{itemize}

4 ¿Cómo puede \textbf{estimarse} desde un \textbf{enfoque frecuentista y
uno bayesiano}, una tasa de \textbf{mortalidad} por causas seleccionadas
en \textbf{distintos estratos de la base}?

~\textbf{Contraargumentos}

\begin{itemize}
\item
  \textbf{Lógica:} La simple estimación de una tasa puede desaprovechar
  las herramientas metodológicas que permiten los enfoques frecuentista
  y bayesiano.
\item
  \textbf{Ética:} Podría suscitarse que al tener resultados ligeramente
  distintos por el cambio de metodología se genere una opinión de valor
  inadecuada en la que se exalte una metodología sobre otra cuando ambas
  son lógicamente válidas pero más convenientes cada una para diferentes
  contextos.
\item
  \textbf{Emocional:} Una población no especializada podría interpretar
  la diferencia de resultados como que una metodología es más adecuada
  que otra.
\end{itemize}

\textbf{Argumentos}

\begin{itemize}
\item
  \textbf{Lógica:} Determinar la tasa de mortalidad brinda una
  comprensión intuitiva de la distribución de la proporción de muerte
  entre causas.
\item
  \textbf{Ética:} Clarificar la concentración de fallecimientos por
  determinadas causas de muerte.
\item
  \textbf{Emocional:} Estimar una tasa de mortalidad para determinadas
  causas permite tener un estadístico que representa las proporciones de
  mortalidad por categoría de forma intuitiva para personas no
  especializadas.
\end{itemize}

De estas cuatro formulaciones, la segunda se considera la más adecuada
como pregunta principal. Esto se debe a que recoge mejor el núcleo del
proyecto: no solamente describir cómo se distribuyen las muertes por
causa, sino estudiar e interpretar cómo cambia la probabilidad de
decremento por categorías de causas a lo largo de los grupos de edad.
Además, mantiene de forma moderada la comparación entre enfoque clásico
y bayesiano, pero sin desplazar el interés principal hacia una discusión
puramente metodológica. Así, la pregunta central sigue estando en el
comportamiento e interpretación de nqx(c) , que es precisamente el
objeto actuarial que se había definido desde el inicio.

\subsubsection{1.1.6 Argumentación a través de los
datos}\label{argumentaciuxf3n-a-travuxe9s-de-los-datos}

Los datos utilizados para este trabajo provienen de la base de
mortalidad de la OMS (En inglés, World Health Organization Mortality
Database). Está dividida en seis partes, cada una recopilando el número
de defunciones clasificadas por rangos de edad y causa de muerte, de
acuerdo al segundo estándar internacional de enfermedades más reciente,
el ICD-10. (International Classification of Diseases). Fuente de
información: WHO Mortality Database. Los datos fueron obtenidos por las
autoridades nacionales de cada país, encargadas de contabilizar y
clasificar las defunciones.

\textbf{Contexto temporal y espacial:} La base de datos registra
anualmente el número de defunciones por país, grupo de edad y causa de
muerte, de acuerdo al ICD-10. Las seis bases unidas contienen
información desde 1950 hasta el año 2023. Para efectos de este trabajo,
se propone trabajar con Centroamérica (Con la aclaración de que se omite
a Honduras del análisis, debido a la falta de registros recientes,
llegando hasta el 2013). En el período 2015-2018.

\textbf{Facilidad de obtener la información:} Las bases de mortalidad de
la OMS se encuentran en su página oficial, junto a un documento que
explica a detalle la dinámica de la misma. Esta guía es necesaria para
interpretar correctamente el código de las causas de muerte y los rangos
de edad en los que están ordenadas las observaciones.

\textbf{Población de estudio:} Las defunciones registradas que forman
parte de los sistemas nacionales obtenidos\textbar{} de la OMS.

\textbf{Muestra Observada:} La cantidad de defunciones registradas,
clasificadas por: Año, país, edad y causa de muerte. Es necesario
aclarar que las defunciones son el agregado de rangos de edades y no por
unidades registradas. Por ende, el proyecto debe de ser interpretado
desde grupos de población.

\textbf{Unidad estadística o individuos:} La unidad estadística es un
grupo poblacional agregado, por ejemplo: ``hombres de 50--54 años, Costa
Rica, 2018, que murieron por una causa determinada''.

\textbf{Descripción de la tabla}

La International Classification of Diseases 10 (ICD-10) es la décima
estandarización de la clasificación de enfermedades propuesta por la
WHO. Esta tiene como finalidad relacionar enfermedades muy complejas en
clasificaciones más amplias y generales (AAPC Thought Leadership Team,
2025). La clasificación de la causa se escribe mediante cuatro
caracteres:

\begin{itemize}
\item
  \textbf{Primer caracter:} Es una letra, esta clasifica de manera más
  amplia la enfermedad.
\item
  \textbf{Segundo caracter:} Es un número. Dentro de la clasificación
  del primer caracter, el segundo caracter desglosa las enfermedades
  asociadas a la clasificación principal.
\item
  \textbf{Tercer caracter:}Es un número. Cumple la misma función que el
  segundo caracter.
\item
  \textbf{Cuarto caracter:} Es un número. Dentro del segundo y tercer
  caracter, se encuentran enfermedades más específicas, las cuales el
  cuarto caracter clasifica. Si la enfermedad no tiene variantes más
  específicas, entonces la enfermedad solo contempla los primeros tres.
\end{itemize}

Por ejemplo: Las letras A y B clasifican ciertas enfermedades
infecciosas y parasitarias. De esta manera se permite agrupar causas de
defunciones específicas en categorías generales:

\begin{longtable}[]{@{}
  >{\raggedright\arraybackslash}p{(\linewidth - 4\tabcolsep) * \real{0.2373}}
  >{\raggedright\arraybackslash}p{(\linewidth - 4\tabcolsep) * \real{0.2627}}
  >{\raggedright\arraybackslash}p{(\linewidth - 4\tabcolsep) * \real{0.5000}}@{}}
\caption{\emph{Ejemplo breve de la clasificación de enfermedades de
acuerdo al ICD-10. Elaboración propia.}}\tabularnewline
\toprule\noalign{}
\begin{minipage}[b]{\linewidth}\raggedright
\textbf{Primer Caracter (Letra)}
\end{minipage} & \begin{minipage}[b]{\linewidth}\raggedright
\textbf{Demás caracteres (Números)}
\end{minipage} & \begin{minipage}[b]{\linewidth}\raggedright
\textbf{Enfermedad}
\end{minipage} \\
\midrule\noalign{}
\endfirsthead
\toprule\noalign{}
\begin{minipage}[b]{\linewidth}\raggedright
\textbf{Primer Caracter (Letra)}
\end{minipage} & \begin{minipage}[b]{\linewidth}\raggedright
\textbf{Demás caracteres (Números)}
\end{minipage} & \begin{minipage}[b]{\linewidth}\raggedright
\textbf{Enfermedad}
\end{minipage} \\
\midrule\noalign{}
\endhead
\bottomrule\noalign{}
\endlastfoot
A-B & 00-99 & Ciertas enfermedades infecciosas y parasitarias \\
A & 00-09 & Enfermedades infecciosas intestinales \\
A & 00 & Cólera \\
A & 001 & Cólera clásica \\
A & 002 & El Tor de cólera \\
A & 009 & Cólera no especificada \\
B & 00-09 & Infecciones virales por lesiones de piel y membrana
mucosa \\
B & 00 & Herpes simple \\
B & 001 & Dermatitis vesicular herpética \\
B & 04 & Viruela del mono (sin variantes específicas) \\
\end{longtable}

Existe una estandarización aún más general para la clasificación de
enfermedades utilizando cuatro dígitos, la cual fue propuesta por
Goerlich (2012), para agrupar de manera más amplia la causa de muerte en
el cálculo de las tablas de vida de decrementos múltiples. Para efectos
de este trabajo, se utilizará este estándar para un mejor manejo de
datos. A continuación una tabla que explica el estándar:

\begin{longtable}[]{@{}
  >{\raggedright\arraybackslash}p{(\linewidth - 4\tabcolsep) * \real{0.1228}}
  >{\raggedright\arraybackslash}p{(\linewidth - 4\tabcolsep) * \real{0.1520}}
  >{\raggedright\arraybackslash}p{(\linewidth - 4\tabcolsep) * \real{0.7251}}@{}}
\caption{\emph{Explicación de la lista de tabulación de mortalidad de la
ICD-10. Obtenido de: (Goerlich, 2012).}}\tabularnewline
\toprule\noalign{}
\begin{minipage}[b]{\linewidth}\raggedright
\textbf{Grupo}
\end{minipage} & \begin{minipage}[b]{\linewidth}\raggedright
\textbf{ICD-10}
\end{minipage} & \begin{minipage}[b]{\linewidth}\raggedright
\textbf{Causa}
\end{minipage} \\
\midrule\noalign{}
\endfirsthead
\toprule\noalign{}
\begin{minipage}[b]{\linewidth}\raggedright
\textbf{Grupo}
\end{minipage} & \begin{minipage}[b]{\linewidth}\raggedright
\textbf{ICD-10}
\end{minipage} & \begin{minipage}[b]{\linewidth}\raggedright
\textbf{Causa}
\end{minipage} \\
\midrule\noalign{}
\endhead
\bottomrule\noalign{}
\endlastfoot
I & A00-B99, R75 & Enfermedades infecciosas y parasitarias \\
II & C00-D48 & Tumores \\
III & D50-D89 & Enfermedades de la sangre y de los órganos
hematopoyéticos, y ciertos trastornos que afectan al mecanismo de la
inmunidad \\
IV & E00-E90 & Enfermedades endocrinas, nutricionales y metabólicas \\
V-VIII & F00-H95 & Trastornos mentales y del comportamiento (V) y
Enfermedades del sistema nervioso y de los órganos de los sentidos
(VI-VIII) \\
IX & I00-I99 & Enfermedades del sistema circulatorio \\
X & J00-J99 & Enfermedades del sistema respiratorio \\
XI & K00-K93 & Enfermedades del sistema digestivo \\
XII & L00-L99 & Enfermedades de la piel y del tejido subcutáneo \\
XIII & M00-M99 & Enfermedades del sistema osteomuscular y del tejido
conjuntivo \\
XIV & N00-N99 & Enfermedades del sistema genitourinario \\
XV & O00-O99 & Embarazo, parto y puerperio \\
XVI & P00-P96 & Afecciones originadas en el periodo perinatal \\
XVII & Q00-Q99 & Malformaciones congénitas, deformidades y anomalías
cromosómicas \\
XVIII & R00-R74, R76-R99, U00-U99 & Síntomas, signos y hallazgos
anormales clínicos y de laboratorio, no clasificados en otra parte. \\
XX & V01-Y89 & Causas externas de mortalidad \\
\end{longtable}

La siguiente tabla explica todo el contenido presente en las columnas de
la base de datos de la WHO:

\begin{longtable}[]{@{}
  >{\raggedright\arraybackslash}p{(\linewidth - 2\tabcolsep) * \real{0.1333}}
  >{\raggedright\arraybackslash}p{(\linewidth - 2\tabcolsep) * \real{0.8667}}@{}}
\caption{\emph{Significado de las columnas de la base de datos. Obtenido
de: World Health Organization (2026).}}\tabularnewline
\toprule\noalign{}
\begin{minipage}[b]{\linewidth}\raggedright
\textbf{Columna}
\end{minipage} & \begin{minipage}[b]{\linewidth}\raggedright
\textbf{Contenido}
\end{minipage} \\
\midrule\noalign{}
\endfirsthead
\toprule\noalign{}
\begin{minipage}[b]{\linewidth}\raggedright
\textbf{Columna}
\end{minipage} & \begin{minipage}[b]{\linewidth}\raggedright
\textbf{Contenido}
\end{minipage} \\
\midrule\noalign{}
\endhead
\bottomrule\noalign{}
\endlastfoot
Country & Código del país. \\
Admin1 & Región específica dentro de cada país (No aplica para
Centroamérica). \\
Subdiv & Tipo de víctimas (Ciudadanos, solo nacionales). No aplica para
Centroamérica. \\
Year & Año del dato registrado. \\
List & Lista de la edición de la ICD utilizada. \\
Cause & Clasificación de la muerte registrada de acuerdo al ICD-10. \\
Sex & 1 para hombre, 2 para mujer, 9 para no especificado. \\
Frmat & Forma de ordenamiento de rangos de edad. \\
IM\_Frmat & Forma de ordenamiento de rangos de edad para menores de un
año. \\
Deaths1 & Muertes para todas las edades. \\
Deaths2 & Muertes a la edad de 0 años. \\
Deaths3 & Muertes a la edad de 1 año. \\
Deaths4 & Muertes a la edad de 2 años. \\
Deaths5 & Muertes a la edad de 3 años. \\
Deaths6 & Muertes a la edad de 4 años. \\
Deaths7 & Muertes en el rango de 5-9 años. \\
Deaths8 & Muertes en el rango de 10-14 años. \\
Deaths9 & Muertes en el rango de 15-19 años. \\
Deaths10 & Muertes en el rango de 20-24 años. \\
Deaths11 & Muertes en el rango de 25-29 años. \\
Deaths12 & Muertes en el rango de 30-34 años. \\
Deaths13 & Muertes en el rango de 35-39 años. \\
Deaths14 & Muertes en el rango de 40-44 años. \\
Deaths15 & Muertes en el rango de 45-49 años. \\
Deaths16 & Muertes en el rango de 50-54 años. \\
Deaths17 & Muertes en el rango de 55-59 años. \\
Deaths18 & Muertes en el rango de 60-64 años. \\
Deaths19 & Muertes en el rango de 65-69 años. \\
Deaths20 & Muertes en el rango de 70-74 años. \\
Deaths21 & Muertes en el rango de 75-79 años. \\
Deaths22 & Muertes en el rango de 80-84 años. \\
Deaths23 & Muertes en el rango de 85-89 años. \\
Deaths24 & Muertes en el rango de 90-94 años. \\
Deaths25 & Muertes a la edad de 95 años o más. \\
Deaths26 & Muertes a una edad no especificada. \\
IM\_deaths1 & Muertes a una edad de 0 días. \\
IM\_deaths2 & Muertes a una edad de 1-6 días. \\
IM\_deaths3 & Muertes a una edad de 7-27 días. \\
IM\_deaths4 & Muertes a una edad de 28-364 días. \\
Pais & Nombre del país de acuerdo a su código. \\
\end{longtable}

A continuación, aclaraciones del cuadro anterior: Para la columna de
Country, en la documentación de la base de datos se encuentra una lista
con los códigos de los países junto a su respectivo nombre. Con la
finalidad de que el trabajo sea eficiente, se presentan sólo los códigos
de Centroamérica:

\begin{itemize}
\tightlist
\item
  Belice: 2045
\item
  Costa Rica: 2140 -El Salvador: 2190
\item
  Guatemala: 2250
\item
  Nicaragua: 2340
\item
  Panamá: 2350
\end{itemize}

Para Centroamérica, las columnas Admin1 y Subdiv vienen vacías (NA).
Esto significa que todas las defunciones registradas no tienen
especificación de región, son datos estandarizados a nivel nacional
(World Health Organization, 2026).

Para las columnas Frmat e IM\_Frmat, si en estas columnas hay un 0 y un
1 respectivamente, los rangos de edad están ordenados de acuerdo a la
tabla anterior. Si posee otro número (Ya sea cualquiera del 1 al 9 para
Frmat, y cualquiera del 2, 8 o 9 para IM\_Frmat), los rangos de edad
presentan un orden distinto. Para efectos de Centroamérica, sólo Costa
Rica posee datos de tipo de orden número 2 en Frmat y 8 para IM\_Frmat
para el año 2020. Por lo tanto, la diferencia sería la siguiente:

\begin{itemize}
\item
  Frmat = 2: No existe la separación de 1,2,3,4 años. Por lo que, el
  agregado de todos los registros de este rango de edad recae en la
  columna de Deaths3 (Pasaría a ser muertes para el rango de edades de 1
  a 4 años). Tampoco existen datos en las columnas Deaths24 y Deaths 25
  (Es decir, 90 años en adelante). Entonces, el agregado de los
  registros para ese rango de edad recae en Deaths 23 (serían muertes
  para la edad de 85 años en adelante).
\item
  IM\_Frmat = 8: No existe la separación de días para muertes en edades
  de menos de un año. Por lo que, cualquier muerte menor a un año recae
  en IM\_deaths1.
\end{itemize}

La columna Pais fue colocada manualmente, para reconocer cada registro
por el nombre del país además de su código.

Finalmente, con respecto a la relación de la tabla con la pregunta
principal de investigación, tenemos que la tabla permite observar, para
cada país, año, grupo de edad y categoría de causa seleccionada, cómo se
comporta la probabilidad estimada de decremento \({}_{n}q_{x}^{(c)}\) y
si ese comportamiento sugiere patrones de aumento, disminución o
estabilidad conforme avanzan los grupos de edad, lo que responde
directamente la pregunta. Además, al incorporar estimaciones obtenidas
desde un enfoque clásico y uno bayesiano, la tabla también hace posible
examinar si ambos marcos ofrecen lecturas semejantes o distintas de esos
patrones.

\subsection{1.2 Revisión
Bibliográfica}\label{revisiuxf3n-bibliogruxe1fica}

\subsubsection{1.2.1 Búsqueda de
Bibliografía}\label{buxfasqueda-de-bibliografuxeda}

Se ha realizado las siguientes combinaciones de palabras clave:

\begin{itemize}
\tightlist
\item
  Base Mortalidad OMS + Causa de Muerte + ICD-10
\item
  Central America + Mortality + Diseases
\item
  Tablas de mortalidad + Bayesiano + Gumpertz, Makeham y Exponencial +
  Latinoamérica
\item
  Probabilidad de decremento + Grupo de edad + Latinoamérica
\item
  Mortalidad por causa de muerte + Cohortes de edad + Latinoamérica
\item
  Estimación bayesiana de mortalidad + Grupo específico de edad + Causa
  de muerte + ICD-10 + Poblaciones pequeñas
\item
  WHO Mortality Database + Bayesian estimation + Age-specific cause of
  death + Central America + small populations
\item
  Specific cause mortality + Gompertz, Makeham and Exponential models +
  Latin America
\item
  Probability of decrement + ICD-10 + Small populations
\item
  Bayesian mortality tables + Age groups + Latin America
\item
  Estimation of specific cause mortality + WHO Mortality Database +
  Exponential models
\end{itemize}

\subsubsection{1.2.2 Construcción de fichas de
literatura}\label{construcciuxf3n-de-fichas-de-literatura}

\textbf{Estimación del riesgo de longevidad mediante distribuciones de
supervivencia transmutadas.}

\begin{itemize}
\item
  \textbf{Título:} Estimación del riesgo de longevidad mediante
  distribuciones de supervivencia transmutadas.
\item
  \textbf{Autor:} Juan Sebastían Orozco Cortés.
\item
  \textbf{Año de publicación:} 2023.
\item
  \textbf{Nombre del tema:} Estimación y proyección de la mortalidad en
  Colombia mediante modelos clásicos y alternativas bayesianas, aplica
  transmutación a los modelos (Descomponer la distribución acumulada
  para reducir los sesgos de la distribución base).
\item
  \textbf{Clasificación:} Metodológica, pues el paper propone un
  procedimiento estructurado paso por paso para la estimación y
  proyección de las tablas de mortalidad mediante modelos como
  Gompertz-Makeham, Perks I y II, exponencial, etc.
\item
  \textbf{Resumen en una oración:} Propone mejorar la estimación de
  mortalidad,con parámetros dinámicos mediante métodos
  clásicos/bayesianos. Argumento central: El paper plantea que la
  mortalidad se puede estimar de manera más precisa y con menos margen
  de sesgo en edades avanzadas mediante distribuciones transmutadas,
  tomando los parámetros de transmutación como series de tiempo.
\item
  \textbf{Problemas con el argumento o el tema:} Si bien se aplican los
  modelos clásicos y bayesianos junto a las distribuciones transmutadas,
  no compara explícitamente la aplicación de los mismos. Además, la
  construcción de los modelos de mortalidad se realiza con observaciones
  anuales, a diferencia de este trabajo, que utiliza observaciones bajo
  grupos de edad y causas de muerte.
\item
  \textbf{Conexión con mi proyecto:} Este paper tiene una relación
  directa con la propuesta planteada en este proyecto, pues muestra que
  las tablas de mortalidad se pueden construir tanto mediante modelos
  clásicos y bayesianos. Esto respalda la idea principal, la cual
  consiste en comparar ambos métodos de inferencia. Añadido a ello, se
  pueden construir las causas de muertes y los grupos de edad para
  efectos de las constantes de los modelos conocidos (Gompertz-Makeham,
  Exp) o parámetros para el caso bayesiano.
\item
  \textbf{Lo que NO dice el autor:} No considera problemas de
  comparabilidad, pues solo construye las tablas de mortalidad para
  Colombia. Además de que sus variables son observaciones individuales y
  no variables agregadas, como los grupos de edad.
\item
  \textbf{Contraste con otra fuente:} En contraste con \emph{The adult
  mortality profile by cause of death in 10 Latin American countries
  (2000--2016)}, el estudio es más de carácter descriptivo, mostrando el
  perfil de mortalidad para adultos mapeado en diez países de
  Latinoamérica y no va más allá de ello, obviando por completo el
  proceso metodológico.
\item
  \textbf{Evaluación de aplicabilidad: 4/5:} La metodología propuesta en
  este documento es aplicable al proyecto, pues involucra los dos tipos
  de inferencia ya mencionados para la estimación de las
  constantes/parámetros. Además del factor de las observaciones
  individuales limita que aplique de la misma forma que en el paper,
  pero se puede adaptar a este trabajo.
\item
  \textbf{Pregunta que le haría al autor:} Precisamente, ¿De qué manera
  adaptaría la propuesta metodológica de su trabajo para la construcción
  de la mortalidad para observaciones agregadas de grupos de edad y
  causas de muerte?
\item
  \textbf{Resumen argumentativo:} El documento muestra que se puede
  mejorar la estimación de la mortalidad en edades avanzadas
  transmutando los modelos de inferencia clásicos y bayesianos. Esta
  diferencia resulta útil para estimar tablas de mortalidad para
  observaciones individuales. Sin embargo, se limita a contextos de
  observaciones agregadas como grupos de edad y a condiciones extras
  como causas de muerte. Por ende, este documento muestra que es posible
  y viable construir las porciones de mortalidad bajo estos modelos,
  argumento clave para la fundamentación de este trabajo. Para efectos
  del mismo, se va a adaptar la metodología propuesta en el paper, pues
  sigue un patrón similar a la idea central.
\end{itemize}

\subsubsection{1.3 Construcción de UVE de
Gowin}\label{construcciuxf3n-de-uve-de-gowin}

\subsubsection{1.4 Parte escrita}\label{parte-escrita}

La pregunta delimitada para el trabajo de investigación constituye ser
\textbf{Análisis del cambio de la probabilidad de decremento}nqx(c)
\textbf{para grupos etarios por~ causas de muerte de acuerdo a la
clasificación ICD-10 en países de Centroamérica {[}2015, 2018{]}.} De
forma directa, se puede afirmar que variables clave para la estimación
de nqx(c)~ constituyen ser:

\(d_x^{(c)}\)~ : Número de vidas saliendo de la población entre las
edades (x,x+1{]}~ debido al decremento c~

\(l_x^{(T)}\)~ : Cantidad de población sobreviviente a la edad exacta x.

\(_nq_x^{(c)}\) : Probabilidad de que una vida de edad (x) deje el grupo
en los siguientes n periodos como resultado de decremento por el factor
(c)

\(_nm_x^{(c)}\) : Tasa central de decremento por la causa (c), el cual
es un promedio ponderado de la fuerza de mortalidad entre las edades
(x,x+1{]}.~

A su vez es conveniente tener presente las siguientes relaciones de
equivalencia que facultará la estimación de la probabilidad de
decremento múltiple; \(_nq_x^{(c)} =d_x^{(c)}/ l_x^{(T)}\). La notación
empleada sigue los estándares actuariales internacionales del material
literario de referencia primario Bowers et al.~(1997).~

Uno de los principales desafíos supondrá la traducción de las
frecuencias de muerte por categorías de causas para diferentes grupos
etarios a una distribución de probabilidad asociada para lo cual se
empleará el uso de la \(_nm_x^{(c)}\), que representa la frecuencia por
unidad de exposición, es decir la intensidad de ocurrencia del evento.
considerando una intensidad constante para los grupos etarios
delimitados en la ICD-10, la acumulación de esa intensidad permite
construir la probabilidad de ocurrencia del evento, y la descomposición
por causa permite obtener la distribución de probabilidad marginal por
causas de decremento.

Dentro de la bibliografía empleada para sustentar las bases del aparato
teórico-metodológico en Dirgová Luptáková y Bilíková (2014) se construye
un modelo de decremento simple de dos estados, a saber; 0:=\{vivo\},
1:=\{muerto\}, donde el estado \{1\} es absorbente pues representa la
única causa de salida del sistema para un individuo de la población,
cabe mencionar que para este modelo en particular se satisface la
propiedad de Markov, por lo que las probabilidades condicionadas de que
el estado futuro \(X_n=y\) solo depende del estado presente
\(X_{n-1}=x\) , lo cual permite ciertas simplificaciones al momento de
realizar cálculos.

Sin embargo, este modelamiento se puede generalizar para más de dos
estado introduciendo de esta forma las probabilidades de decremento
múltiples, con esta última idea es que se persigue articular el proyecto
de investigación propuesto, para el cual se comprenderá un modelo de
diversas causas de salidas c que corresponderá a los estados absorbentes
que serán las diversas causas de muerte basados en la nomenclatura de la
ICD-10.

Cabe indicar que realizando ciertos supuestos tales como que al
condicionar el estado actual del agregado de la población de estudio y
al grupo etario al que pertenece actualmente, el futuro no depende de
cómo llegó ahí, es decir que se satisfaga la propiedad de Makov, se
podría implementar las facilidades para ejecutar cálculos siempre y
cuando estos correspondan asunciones sensatas dentro del modelo.~

Por otro lado el establecimiento de un modelo multiestado para
representar las probabilidades de decremento viene motivada porque según
afirma Goerlich Gisbert (2012) a ``estudios que tratan de ajustar la
esperanza de vida por estados de salud están muy limitados por la escasa
disponibilidad de información estadística~ y las dificultades de
comparación intertemporal'' (p.~6) , está crítica está fundamentada
esencialmente en que la esperanza de vida es una medida dispersión muy
asimétrica, aunado a su vez que la \(E[T_{x}]\)es sencillamente la edad
media a la muerte, es decir cuando ocurre y no aporta información
detallada que especifique las causas que ocasionan el perecimiento ni
mucho menos el cómo cambian estas.

En este sentido, extendiendo las tablas de mortalidad de decrementos
únicos a tablas por decrementos múltiples, es decir la representación de
las probabilidades de transición a un estado absorbente \{y\} por
múltiples causas comprendidas en la ICD-10, se permite comprender cómo
han variado las probabilidades de mortalidad entre grupos etarios al
considerar diferentes patologías de base de la cohorte de estudio; para
este fin se seguirán los lineamientos metodológicos sugeridos por
Goerlich Gisbert en Tablas de vida de decrementos múltiples: Mortalidad
por causas en España (1975--2008) (Goerlich Gisbert, 2012), pero
enfocándonos en una perspectiva Centroamericana.

\newpage

\section{Anexos}\label{anexos}

\subsection{Anexo 1: Preparación}\label{anexo-1-preparaciuxf3n}

\subsubsection{Log de contribuciones}\label{log-de-contribuciones}

\begin{verbatim}
Warning: package 'kableExtra' was built under R version 4.5.3
\end{verbatim}

\begin{longtable*}[t]{>{\raggedright\arraybackslash}p{5cm}>{\raggedright\arraybackslash}p{10cm}}
\toprule
Integrante & Contribución\\
\midrule
Kevin David Calderón Martínez & Argumentación a través de datos. 1.1.6. Búsqueda de bibliografía 2.1.1. Fichas bibliográficas 2.1.2: Estimación del riesgo de longevidad mediante distribuciones de supervivencia transmutadas y The adult mortality profile by cause of death in 10 Latin American countries. Log de contribuciones (Anexo) y Diario de aprendizaje 5\\
Gabriel de Jesús Chaves Esquivel & Identificación de tensiones 1.1.3. UVE de Gowin 1.3. Búsqueda de bibliografía 2.1.1 . Audiencia del proyecto (Anexo). Log de contribuciones (Anexo)\\
Sebastián Miranda Ramírez & Definición de la idea 1.1.1. Conceptualización de la idea 1.1.2. Reformulación de la idea en modo pregunta 1.1.4. Búsqueda de bibliografía 2.1.1 . Registro de uso de IA (Anexo). Log de contribuciones (Anexo). Búsqueda de papers y material al respecto de las mortalidades de decrementos múltiples.\\
Benjamín Gutiérrez Padua & Pregunta de investigación 1. Argumentación de la pregunta 1.1.5. Búsqueda de bibliografía 2.1.1. Fichas bibliográficas 2.1.2: WHO Mortability Database. Parte de escritura 3. Log de contribuciones (Anexo)\\
\bottomrule
\end{longtable*}




\end{document}
